\section{Normalizing the exceptional points in a patching chain}

Let $x$ be a point of the original ordinary determinant problem,
$A=k[[t]]$ its curve DVR and $E=k((t))$. Write $G'$ for the
distinguished subgroup in the source and $P$ for the invertible
symmetric form of $x$. Assume the original hypotheses of
$\mathrm{Norm}$, in particular $n>0$, $\ell\ne2$ and
$I_u\subseteq r.I_p(u)$ for each prescribed inertia group.
CIP10 treats the absolutely irreducible case. The other allowed
case has two absolutely irreducible constituents on $G'$. To use
the point in a chain we need one integral basis respecting those two
blocks and the polarization, as well as the original ordered flags.

% BEGIN N69E1
\paragraph{N69E1. The image and the polarization have the same two blocks.}
For a type-E point, choose the splitting of the source over $E$,
with block sizes $a,b>0$, $a+b=n$, and reorder the coordinates so
the first block comes first. On $G'$ the two representations are
absolutely simple and have no nonzero intertwiner from the first to
the second. Over $\bar E$ they are nonisomorphic. The image algebra
of $\bar E[G']$ in $M_a(\bar E)\times M_b(\bar E)$ is the entire
product. Indeed each projection is onto by the density theorem for a
simple module over an algebraically closed field. Their annihilators
in $\bar E[G']$ are distinct maximal two-sided ideals, since a full
matrix algebra has just one simple module up to isomorphism.
The Chinese remainder theorem for these comaximal ideals gives
surjectivity onto the product. Extension of scalars preserves the
rank of the span, so the $E$-span is also the full block algebra.

Transport $P$ by inverse-transpose and inverse of the splitting
matrix. Each block has its own nonsingular polarization with the
same prescribed multiplier. Comparing either off-diagonal block
of the transported pairing with the constituent's own pairing
therefore gives an intertwiner between the two constituents on
the conjugate of $G'$. Conjugation preserves $G'$ and is bijective
there. The two constituents have no such intertwiner in either
direction by simplicity and nonisomorphism. Hence the cross terms
are zero and each diagonal block form is nonsingular and symmetric.
The decomposition was given on the full representation group, which
ensures that the two summands themselves are stable there.
% END N69E1

% BEGIN N69E2
\paragraph{N69E2. An invariant unimodular lattice in each block.}
Project the original stable lattice onto each summand of the
splitting in N69E1. These images are finitely generated, span
their summands over $E$, and are stable under the group and its
inverses. Fix such a lattice $D$ and its block pairing $H$.
Its dual $D^\vee$ is stable too. Indeed the conjugation in the
polarization permutes the group elements, the multiplier takes
unit values, and $D$ is stable under every inverse. Choose an
orthogonal $A$-basis of $D$ for $H$, with diagonal entries
$t^{d_i}v_i$, $d_i\in\mathbb Z$, $v_i\in A^\times$. To see this,
scale the form until its least entry valuation is zero. A unit
diagonal entry gives a primitive pivot. If only an off-diagonal
entry is a unit, the sum of the two basis vectors has unit
self-pairing because $2$ is invertible. Split off the resulting
orthogonal line, and repeat with the restricted form. Undo the
scalings to obtain the claimed exponents.

The dual lattice has generators $t^{-d_i}e_i$. For an action
matrix $g=(g_{ij})$, preservation of both lattices gives
\[
       v_t(g_{ij})\ \ge\ \max(0,d_j-d_i).
\]
Extend $A\to k[[u]]$ by $t\mapsto u^2$ and take the lattice with
basis $u^{-d_i}e_i$. Its $ij$ matrix entry has valuation at least
$2\max(0,d_j-d_i)+d_i-d_j=|d_j-d_i|$. This applies to $g^{-1}$
too. Its Gram entries on the diagonal are the units $v_i(u^2)$.
Enlarge the finite residue field once to contain square roots of
all their residues in both blocks. Hensel lifting and unit rescaling
in each block give Gram matrix $1$. Thus one finite injective local
map $j:A\to A'=k'[[u]]$ gives integral invertible matrices of both
blocks, with no off-diagonal entries. This proves the lattice and
pairing assertions, including the original multiplier. Entrywise
adic continuity follows by conjugating by a fixed fraction-field
matrix: the valuations of its finitely many entries and those of
its inverse bound the loss of precision uniformly in $g$.
% END N69E2

% BEGIN N69E3
\paragraph{N69E3. All the fields of the original normalization.}
After this conjugation write the representation as $y.\rho$.
In either block the span statement of N69E1 still holds over
$\operatorname{Frac}(A')$, and together the blocks still span
their product. Multiplying an integral block matrix by a sufficiently
large power of $u$ clears the finitely many denominators in an
expression in group elements of $G'$. This gives the required
integral span. Take $y.\psi_{v,i}=j\psi_{v,i}$. Their distinctness
and the exclusions $\psi_{v,i}\ne\psi_{v,j}\varepsilon$ for $i<j$
are preserved by injectivity. The prescribed torsion subgroup
has trivial image under each point character by the original
torsion condition and its reduction in characteristic $\ell$
(as in CIP8). If $g\in I_u$, the original inertia condition kills
$g$. Conjugation carries this identity to the new matrices.

At each $v$, transport the original ordered fraction-field flag.
Its factors are the distinct characters $j\psi_{v,i}$. The two block
modules have disjoint sets of composition factors. Any submodule
of their direct sum is therefore the sum of its intersections with
the two blocks: the projections modulo those intersections are
isomorphic subquotients of the two modules, hence zero. Apply this
to each term of the ordered flag. Its dimension increases in
exactly one block at each step. Intersect the flags with the two
integral lattices. Each intersection is saturated, and successive
rank-one quotients in the growing block are free over the DVR.
Extending bases successively gives an integral basis in the
specified order, each column supported in its indicated block.
Its diagonal characters are the given ones since they agree over
the fraction field.

For every group element the characteristic polynomial of $y.\rho$
is the image of the old one, by conjugacy and the injection of the
DVR into its fraction field. The determinant character-coefficient
criterion for group determinants gives equality of the full
polynomial laws. Equivalently, for a finite list of group elements,
compare the determinants of the two generic linear combinations
over $A'[X_1,\ldots,X_s]$ by conjugating over its fraction field.
Injectivity puts the polynomial identity back in the polynomial
ring. Specializing the variables gives the law over every
coefficient algebra. Separability transfers under field extension,
so the original open-set test $\mathrm{N3}$ transfers. For
$\mathrm{N4}$ retain the \emph{old} subgroup $F.D_0$.
The equations $\varepsilon(g)\rho(g)Z=Z\rho(g)$ form a system
of homogeneous linear equations in $n^2$ variables. A finite
subset of their rows spans the row space over $E$. Its kernel
is zero by $x.\mathrm{N4}$ and stays zero after field extension.
Conjugate and restrict back to integral matrices to get
$y.\mathrm{N4}(F.D_0)$. This supplies the native normalized point
on $\mathrm{reIu}(F,I)$, with $a$ its first block size. In type I,
CIP10 supplies $a=n$. The two cases exhaust the source's
$\mathrm{NiceI}\lor\mathrm{NiceE}$ hypothesis, proving
$\mathrm{Norm}$ at every allowed rank.
% END N69E3

% BEGIN N69E4
\paragraph{N69E4. Patching a chain without changing any of its prime kernels.}
Retain the source inputs $\mathrm{Choice}$, the
first-order dual bound and Taylor--Wiles obstruction bound, and
$\mathrm{Raw}$. We show how these give the full
$\mathrm{PfromAut}$ with N69E3. Fix its level $V$, its $B,J$,
the cohomological and local reciprocity data, the nontrivial residual
cyclotomic test, and $\mathrm{HasAut}(B,J)$. Given any finite list
$x_0,\ldots,x_h$ of nice points of either kind and their maps $f_i$
out of $B$, first use N69E3 with the inertia groups of the
arithmetic frame. They agree with the required original $r.I_p$
there. The arithmetic place theorem supplies the finiteness
condition $\mathrm{PhiCond}$ by CAP13. Apply $\mathrm{Choice}$ once to
the entire normalized list. It gives one list of detecting elements,
integral distinct ordered roots at every chosen element for each
point, injective finite local coefficient extensions and equal
residue-field cardinalities for the extended points. The arithmetic
place theorem now returns \emph{one} Taylor--Wiles system with those
elements. Use $\mathrm{HasAut}(B,J)$ on this system and its actual
places to get its automorphic module.

Let $z_i$ denote the final points and $j_i$ the two successive
coefficient maps composed. The determinant and every ordered
character map to those of $z_i$ under $j_i$. The extended
representation kills the relative pro-$\ell$ residual required
in the two bounds. Indeed on the subgroup $H$ cut out by the
residual characters and cyclotomic condition the specialization
at the constant coefficient has characteristic polynomial
$(X-1)^n$ on every element, by the residual determinant identity.
Locality of $j_i$ makes constant coefficient commute with the
induced embedding of residue fields. Thus reduction of the image
of $H$ is a finite unipotent group in characteristic $\ell$
and is an $\ell$-group. Kernels of successive integral congruence
reductions are $\ell$-groups too. All finite congruence quotients
give open kernels in the index-two representation group. Take
their normal cores in the full frame group. The image of $H$ in
the resulting finite quotient still embeds in a product of
$\ell$-groups, since $H$ is normal. These normal subgroups are
among the tests defining the relative pro-$\ell$ residual.
Separatedness now gives the assertion.
The first-order bound now supplies $B1$ with parameter
$n[L^+:\mathbb Q]+nq$, for $q$ the number of detecting elements.
The second supplies $B2$ for the very same point and auxiliary data.
The maps of the resulting patched points are $j_i f_i$ and have
kernel \emph{exactly} $\ker f_i$ by injectivity.

The adaptedness required by $\mathrm{Raw}$ follows from the
description of $H$ and the arithmetic inertia comparison.
The original universal determinant and character coefficients
topologically generate $B$: modulo each maximal-ideal power every
element belongs to the image of their generated subalgebra.
They are among the pair values used by $\mathrm{Raw}$.
Its chain implication therefore applies to these original kernels
and the original ideal $J$.
% END N69E4

% BEGIN N69E5
\paragraph{N69E5. The two prime tests supplied by the patching chain.}
For a type-I normalized point its full matrix span over the
fraction field gives $n^2$ group matrices forming a basis.
The trace pairing on $M_n$ is nondegenerate, since matrix units
pair with their transposes. Hence the determinant of their trace
Gram matrix is nonzero at this point. The corresponding universal
element avoids every prime beneath its kernel. This is
$\mathrm{DiscSt}$ for each such prime. Use N69E4 with one nice
type-I point and no edges. When $J\subseteq\ker f$, the prime
$\ker f$ itself supplies the starting discriminant prime. Every
prime underneath passes the terminal discriminant test just proved.
The chain conclusion puts $J$ in all of them, exactly
$\mathrm{P1core}(J)$.

For a list of nice exceptional points use N69E4 at its full length.
A stable prime is a discriminant prime by the same trace-pairing
criterion over the algebraic closure of its residue field: the
semisimple realization is absolutely irreducible exactly when its
image spans the full matrix algebra. A thick prime gives the
required failure of semisimplicity of the specialized group-algebra
quotient in the definition of $\mathrm{NonSS}$, using the map to
that same algebraic closure with kernel the specified prime.
The augmentation ideals before and after the arithmetic change of
frame agree because its prescribed inertia groups agree and the
determinant and ordered-character relations are the same.
Thus each connecting thick prime lies below both of the original
adjacent kernels and is an allowed edge. Use the given stable
prime containing $J$ below the first kernel as the starting prime.
Every stable or thick prime below the last kernel passes the
terminal test. This gives $J\subseteq\mathfrak q$ for each such
prime, exactly $\mathrm{P3core}(J)$, also when $h=0$.
Together the two arguments prove the all-rank implication with the
four patching inputs stated in N69E4. In particular HIA3 can use
this implication with the Hida ideal supplied by HIA2.
% END N69E5
